% !TEX root = ../../Main.tex
\section{Contributions}

The contributions of this work are the following.

\begin{itemize}
\item \textbf{A trading framework where one strategy implementation runs in three places.} The system has two parts. A Python library holds the strategy, the models and the market abstractions. A C\# component runs inside cTrader \cite{noauthor_forex_nodate} and connects cTrader to the library through shared memory. The same strategy code runs without change in an offline backtest over stored market data, in the backtesting tab of cTrader and on a live account.

\item \textbf{A market database built from quote-level data.} The historical data was collected from cTrader through the same connector and stored locally. For the seven pairs studied here, it holds about 1.66 billion quotes, which take 295 GB. It also holds 33.0 million bars: 32.5 million at one minute, 547,835 at one hour, 23,386 daily and 1,050 monthly. All of them cover January 2014 to August 2026. For each run, the engine uses the finest series available. The price path inside a bar is therefore rebuilt from the quotes themselves and is not interpolated.

\item \textbf{A backtesting engine that charges what the account charges.} The engine moves forward tick by tick instead of from one bar close to the next. The spread is the one quoted at the moment of the trade. It is read from the recorded quotes and is not assumed. The commission is computed from the traded volume at the rate of a raw-spread retail account. Stop and target levels are checked against the price path inside the bar instead of at the bar close. This reproduces the full cost structure of a swap-free raw-spread account, which is a type of account that retail brokers really offer. On such an account, the trader pays two charges: the spread and the commission. Both are applied here at the rates that this account pays. To our knowledge, no published work that applies deep reinforcement learning to currency trading combines quote-level data, a measured spread and a real commission schedule in this way. \Cref{Tables:BasicRelatedWork} and \Cref{Tables:AdvancedRelatedWork} summarise what the reviewed studies report. Charging these costs also means that the returns reported in \Cref{Chapter:ExperimentalResults} cannot be compared directly with results obtained on bar closes with a fixed spread or no spread. For the same policy, the returns here are lower than those results would be.

\item \textbf{DDPG applied to the seven major currency pairs under one protocol.} One agent was trained for each pair, with the same architecture, observation design and reward. Each agent was evaluated from January 2015 to January 2026 with the same cost model. The pairs can therefore be compared with each other. The design is tested for generality and is not tuned to one instrument.

\item \textbf{Positive and risk-controlled results on all seven pairs.} Over the eleven years, the total return was between $+15.7\%$ on USD/JPY and $+226.6\%$ on USD/CAD. Every pair stayed profitable when the evaluation was repeated with five different starting balances. The best pair on a risk-adjusted basis was USD/CHF. It had a Sharpe ratio \cite{sharpe_sharpe_1994} of $0.80$, a Sortino ratio \cite{sortino_performance_1994} of $1.55$, a Calmar ratio \cite{bacon_practical_2023} of $0.49$ and a maximum drawdown of $8.2\%$ over 220 trades.

\item \textbf{An attribution of the returns, including one negative result.} Each model's yearly return is regressed on the yearly movement of its own currency pair. This separates the return that comes from the policy from the return that comes from exposure to the market. Six of the seven pairs had a positive intercept, and three of those had a market sensitivity close to zero. The exception is USD/JPY. Its model follows the underlying pair almost exactly and adds nothing to it. This case is reported here, together with the mechanism that causes it.
\end{itemize}